\documentclass[12pt]{article}
\usepackage[margin=0.6in,top=0.85in,bottom=0.55in]{geometry}
\usepackage{lmodern}
\usepackage{microtype}
\usepackage{fancyhdr}
\usepackage{titlesec}
\usepackage{enumitem}
\usepackage{lastpage}
\usepackage[hidelinks]{hyperref}

\setlength{\parindent}{0pt}
\setlength{\parskip}{0.25em}
\setlength{\headheight}{26pt}

\titleformat{\section}{\large\bfseries\scshape}{}{0pt}{}

\pagestyle{fancy}
\fancyhf{}
\fancyhead[C]{%
\footnotesize
Student Name: \hspace{2.3in} \quad Assignment ID: U1L2LL \quad Page \thepage\ of \pageref{LastPage}
}
\renewcommand{\headrulewidth}{0.5pt}
\renewcommand{\footrulewidth}{0pt}

\newcommand{\answerbox}[2]{%
\vspace{0.15em}
\noindent\textbf{#1}\par
\vspace{0.05em}
\noindent\fbox{%
\begin{minipage}[t][#2]{\dimexpr\textwidth-2\fboxsep-2\fboxrule}
\rule{\textwidth}{0pt}
\vspace{#2}
\end{minipage}
}\par
\vspace{0.25em}
}

\begin{document}

\begin{center}
{\LARGE\bfseries Week 2 Logic Lab}\par\vspace{0.05em}
{\large\scshape Terms, Ambiguity, and Equivocation}\par\vspace{0.1em}
\rule{0.78\textwidth}{0.5pt}\par\vspace{0.15em}
\end{center}

\noindent\textbf{Purpose:} Practice Week 2 skills on \textbf{new} examples (adapted from standard logic-text and classroom cases). Find when a word carries more than one idea, name the senses, spot equivocation, and \textbf{fix the term} before judging whether a conclusion follows. Do not reuse the bank, liberal, equal, or tolerance examples from the reading handout.

\vspace{0.3em}

\section*{Part 1: Two Meanings in Context}

\textbf{Part 1 Q1:} In this sentence, \textbf{right} is used twice with different meanings: ``You have a \textbf{right} to speak at the meeting, but that does not mean you are \textbf{right} about the policy.'' State both meanings of \textbf{right}.

\answerbox{Part 1 Q1 ANSWER BOX}{1.55in}

\textbf{Part 1 Q2:} In this sentence, \textbf{free} is used twice with different meanings: ``In a \textbf{free} country, admission to the public museum on holidays should be \textbf{free}.'' State both meanings of \textbf{free}.

\answerbox{Part 1 Q2 ANSWER BOX}{1.65in}

\newpage

\section*{Part 2: Spot the Equivocation}

\textbf{Part 2 Q1:} This argument is a standard textbook case (light as weight vs.\ light as brightness): ``What is light cannot be dark. Feathers are light. Therefore feathers cannot be dark.'' Identify the shifting term, both meanings, and explain why the conclusion does not follow.

\answerbox{Part 2 Q1 ANSWER BOX}{2.45in}

\textbf{Part 2 Q2:} This argument is a standard classroom case (star as celestial body vs.\ star as celebrity): ``All stars are balls of gas in space. The lead singer is a star. Therefore the lead singer is a ball of gas in space.'' Identify the shifting term, both meanings, and explain why the argument fails.

\answerbox{Part 2 Q2 ANSWER BOX}{2.45in}

\newpage

\section*{Part 3: Fix the Term}

\textbf{Part 3 Q1:} This argument is a classic philosophy-text example (end as purpose vs.\ end as termination): ``The end of a thing is its perfection. Death is the end of life. Therefore death is the perfection of life.'' For the word \textbf{end}, state one clear meaning for each use. Then say whether the conclusion follows once the meanings are fixed.

\answerbox{Part 3 Q1 ANSWER BOX}{2.55in}

\newpage

\section*{Part 4: Argument in a Paragraph}

\textbf{Part 4 Q1:} Read the paragraph (adapted from a standard ``nobody'' equivocation case). Find where \textbf{nobody} shifts meaning, name both senses, and explain why the final claim does not follow reliably.

\vspace{0.2em}

\noindent\fbox{%
\begin{minipage}{\dimexpr\textwidth-2\fboxsep-2\fboxrule}
\vspace{0.35em}
After the student council election, Jordan said the loss made him feel like a nobody. The next day, a campaign flyer stated that nobody is perfect, so Jordan must be perfect in some respect. The yearbook editor praised the flyer for proving Jordan's excellence. Jordan replied that the flyer only proved someone had confused two different meanings of one word.
\vspace{0.35em}
\end{minipage}
}

\answerbox{Part 4 Q1 ANSWER BOX}{2.75in}

\newpage

\section*{Part 5: Review — Week 1 Four-Step Test}

\textbf{Part 5 Q1:} ``The mayor has served in office longer than any other council member, so the mayor's new zoning plan must be lawful.'' Use the four-step test: conclusion, stated reason, hidden assumption, follows.

\answerbox{Part 5 Q1 ANSWER BOX}{2.35in}

\textbf{Part 5 Q2:} ``This clinic received the highest patient satisfaction scores, so its medical advice must be scientifically correct.'' Use the four-step test: conclusion, stated reason, hidden assumption, follows.

\answerbox{Part 5 Q2 ANSWER BOX}{2.35in}

\newpage

\textbf{Part 5 Q3:} Read the paragraph. Ignore distractors. Parse the main argument with the four-step test (conclusion, stated reason, hidden assumption, follows).

\vspace{0.2em}

\noindent\fbox{%
\begin{minipage}{\dimexpr\textwidth-2\fboxsep-2\fboxrule}
\vspace{0.35em}
The hiring panel ate lunch in the courtyard while interns rearranged chairs. One member said her cousin dislikes early interviews. The panel chair noted that the finalist had answered every technical question correctly and that no other candidate had completed the same skills trial. A board observer posted a meme about trick questions online. The panel's official note said the finalist should be hired because proven test performance guarantees long-term success in the role.
\vspace{0.35em}
\end{minipage}
}

\answerbox{Part 5 Q3 ANSWER BOX}{2.85in}

\end{document}